\section{Arzelà--Ascoli判据}
\label{sec:fcover-arzela-ascoli}

上一节的帽形函数说明，函数类缺少共同的连续性尺度时，有限个位置上的值无法代表整个函数。本节把这一观察变成判据：在紧度量空间$K$上，连续函数类在一致距离下全有界，当且仅当它一致有界并且等度连续。全节中$C(K)$表示$K$上实值或复值连续函数的全体，距离取$\lVert f-g\rVert_\infty$；两种取值情形的论证相同，只在量化取值时分别处理实轴上的区间与复平面上的正方形。

\paragraph{一致有界与等度连续}

覆盖要求有限个代表照顾到每个成员，因此类中函数至少要共同受到两方面的约束：取值不能任意大，变化不能任意快。第一个约束是幅度上的。

\begin{definition}[一致有界]
\label{def:fcover-uniformly-bounded}
设$K$是集合，$\mathcal F$是$K$上一族实值或复值函数。若存在常数$M$，使$\lvert f(x)\rvert\leq M$对所有$f\in\mathcal F$和所有$x\in K$成立，就称$\mathcal F$\term{一致有界}（uniformly bounded）。
\end{definition}

对$\mathcal F\subseteq C(K)$，一致有界就是$\sup_{f\in\mathcal F}\lVert f\rVert_\infty<\infty$，即$\mathcal F$含于一个一致范数球中。上一节的帽形函数说明，仅有这一条还不够。

第二个约束是变化上的，它就是第~\ref{chap:metric-space-maps}章的等度连续性：对每个$x\in K$和$\varepsilon>0$，有一个对全族都有效的半径$\delta$，使$d(x,y)<\delta$时$\lvert f(x)-f(y)\rvert<\varepsilon$对所有$f\in\mathcal F$成立。那里还证明了，定义域紧时这个半径可以进一步取得与$x$无关。下面使用的正是这一统一形式：对每个$\varepsilon>0$，存在$\delta>0$，使
\begin{equation}
d(x,y)<\delta\ \Longrightarrow\ \lvert f(x)-f(y)\rvert<\varepsilon\qquad\text{对所有}\ f\in\mathcal F\ \text{及}\ x,y\in K.
\label{eq:fcover-uniform-equicontinuity}
\end{equation}

最常见的等度连续族共用一个Lipschitz常数。若每个$f\in\mathcal F$都满足$\lvert f(x)-f(y)\rvert\leq L\,d(x,y)$，取$\delta=\varepsilon/L$即可满足式\eqref{eq:fcover-uniform-equicontinuity}。反过来，上一节的帽形函数$\varphi_k$不等度连续：对任意$\delta>0$，取$k$使$r_k<\delta$，点$x=c_k$与$y=c_k+r_k$相距$r_k<\delta$，却有$\lvert\varphi_k(x)-\varphi_k(y)\rvert=1$。$[0,1]$上的$x^n$也是这样，第~\ref{chap:functions-as-vectors}章已看到它们一致有界却没有一致收敛的子列。

\paragraph{判据与证明}

两个约束合起来恰好刻画全有界性。在陈述之前回忆两件事：$C(K)$在一致距离下完备（第~\ref{chap:normed-spaces}章）；在完备度量空间中，子集相对紧（闭包紧）当且仅当它全有界（第~\ref{chap:metric-spaces}章）。

\begin{theorem}[Arzelà--Ascoli判据]
\label{thm:fcover-arzela-ascoli}
设$(K,d)$是非空紧度量空间，$\mathcal F\subseteq C(K)$，距离取一致距离。则下列三条等价：
\begin{enumerate}
\item $\mathcal F$一致有界且等度连续；
\item $\mathcal F$全有界，即对每个$\varepsilon>0$都有$N_{\mathcal F}(\varepsilon,\mathcal F)<\infty$；
\item $\mathcal F$在$C(K)$中相对紧。
\end{enumerate}
进一步，$\mathcal F$在$C(K)$中紧，当且仅当它满足上述条件并且是$C(K)$的闭子集。
\end{theorem}

证明的核心是有限化：等度连续允许只在有限个位置上观察函数，一致有界允许把这些位置上的值归入有限个小块，于是整类函数只留下有限种“记录”。

\begin{proof}
由1推出2。设$\lvert f\rvert\leq M$对所有$f\in\mathcal F$成立，给定$\varepsilon>0$。由式\eqref{eq:fcover-uniform-equicontinuity}取$\delta>0$，使$d(x,y)<\delta$时$\lvert f(x)-f(y)\rvert<\varepsilon/4$对所有$f$成立。紧空间全有界，所以有有限个点$x_1,\ldots,x_p\in K$，使每个$x\in K$都满足某个$d(x,x_i)<\delta$。

再把取值集合$\{z:\lvert z\rvert\leq M\}$分成有限个直径不超过$\varepsilon/4$的小块：实值时用长度不超过$\varepsilon/4$的小区间，复值时用边长不超过$\varepsilon/(4\sqrt2)$的小正方形，并约定各块互不相交。对每个$f\in\mathcal F$，记下$f(x_1),\ldots,f(x_p)$分别落在哪一块，这一串块编号称为$f$的记录。块数有限、位置数有限，所以记录只有有限种。

对每一种实际出现的记录，从$\mathcal F$中任取一个具有该记录的成员作为代表。若$f$与代表$g$记录相同，则$\lvert f(x_i)-g(x_i)\rvert\leq\varepsilon/4$对每个$i$成立。对任意$x\in K$，取$x_i$使$d(x,x_i)<\delta$，由三角不等式，
\[
\lvert f(x)-g(x)\rvert\leq\lvert f(x)-f(x_i)\rvert+\lvert f(x_i)-g(x_i)\rvert+\lvert g(x_i)-g(x)\rvert<\frac{3\varepsilon}{4}.
\]
所以有限个代表构成中心在$\mathcal F$中的$\varepsilon$覆盖，$N_{\mathcal F}(\varepsilon,\mathcal F)<\infty$。

由2推出1。取一个有限的$1$覆盖$f_1,\ldots,f_N\in\mathcal F$，每个$f$都满足某个$\lVert f-f_j\rVert_\infty\leq1$，故$\lVert f\rVert_\infty\leq\max_j\lVert f_j\rVert_\infty+1$，这就是一致有界。为证等度连续，给定$\varepsilon>0$，取有限的$\varepsilon/3$覆盖$f_1,\ldots,f_N\in\mathcal F$。紧集上的连续函数一致连续（第~\ref{chap:metric-space-maps}章），有限个函数可以共用一个$\delta$，使$d(x,y)<\delta$时$\lvert f_j(x)-f_j(y)\rvert<\varepsilon/3$对每个$j$成立。对任意$f$，选$f_j$使$\lVert f-f_j\rVert_\infty\leq\varepsilon/3$，则
\[
\lvert f(x)-f(y)\rvert\leq\lvert f(x)-f_j(x)\rvert+\lvert f_j(x)-f_j(y)\rvert+\lvert f_j(y)-f(y)\rvert<\varepsilon .
\]

2与3等价。$C(K)$完备，由前面回忆的相对紧判据即得。

最后，紧集是闭集；反过来，若$\mathcal F$相对紧且闭，则$\mathcal F$等于它的闭包，因而紧。
\end{proof}

两个方向的证明各说明了一件事。充分性的构造是可以计数的：位置数$p$由定义域在尺度$\delta$下的覆盖数决定，块数由幅度$M$与精度决定，记录种数至多是块数的$p$次方。第~\ref{sec:fcover-construction}节将在$[0,1]^d$上把这个计数写成显式公式。必要性则说明有限覆盖必然带来共同的连续性尺度，因为有限个函数总能共用一个$\delta$，一致距离再把这个$\delta$传给整个类。

\paragraph{相对紧性与闭性}

定理的最后一句把紧性分成两部分：一致有界与等度连续保证闭包紧，函数类自身紧还需要闭性。这两部分确实不同，下面各举一例。

先看闭的情形。设$L,M>0$，令$\mathcal F_{L,M}$为$K$上满足$\lvert f\rvert\leq M$且$\lvert f(x)-f(y)\rvert\leq L\,d(x,y)$的全体函数。它一致有界且等度连续。若$f_n\in\mathcal F_{L,M}$一致收敛到$f$，则对每一对$x,y$，不等式$\lvert f_n(x)\rvert\leq M$与$\lvert f_n(x)-f_n(y)\rvert\leq L\,d(x,y)$在$n\to\infty$时保持，所以$f\in\mathcal F_{L,M}$。因此$\mathcal F_{L,M}$闭，从而在$C(K)$中紧。

再看不闭的情形。在$[0,1]$上取满足$\lvert f\rvert\leq1$、$\lvert f'\rvert\leq1$的全体连续可微函数。由中值定理，它们共用Lipschitz常数$1$，所以这一类相对紧。但它不闭：令$g_n(x)=\sqrt{(x-1/2)^2+1/n^2}$（$n\geq2$），则$g_n$连续可微，$\lvert g_n'(x)\rvert=\lvert x-1/2\rvert/g_n(x)<1$，且$g_n\leq\sqrt{1/4+1/4}<1$。又
\[
0\leq g_n(x)-\Bigl\lvert x-\frac12\Bigr\rvert=\frac{1/n^2}{g_n(x)+\lvert x-1/2\rvert}\leq\frac1n,
\]
所以$g_n$一致收敛到在$1/2$处不可导的$\lvert x-1/2\rvert$。这一类的闭包是紧的，它自身却不紧，极限跑到了类外。

\paragraph{定义域紧性的作用}

判据中的紧定义域同样不能省略。证明充分性时，我们用有限个点$x_1,\ldots,x_p$控制了整个定义域，这一步依赖$K$的全有界性。若定义域是$\mathbb R$，即使函数类一致有界且共用一个Lipschitz常数，也可能没有有限覆盖。

取$\varphi(x)=\max\{0,1-\lvert x\rvert\}$及平移族$\varphi_n(x)=\varphi(x-3n)$，$n\in\mathbb Z$。全族满足$0\leq\varphi_n\leq1$，Lipschitz常数都是$1$。但$\varphi_n$与$\varphi_m$（$n\neq m$）的支撑相距至少$1$，在$3n$处两者相差$1$，所以$\sup_x\lvert\varphi_n(x)-\varphi_m(x)\rvert=1$。半径小于$1/2$的球至多含一个成员，覆盖数为$+\infty$。函数之间的差异跑向了无穷远，有限个位置看不到它。
